\documentclass[12pt,a4paper]{article}

\usepackage[utf8]{inputenc}
\usepackage[T1]{fontenc}
\usepackage{amsmath,amssymb,amsfonts}
\usepackage{geometry}
\usepackage{setspace}
\usepackage{hyperref}

\geometry{margin=2.5cm}
\onehalfspacing

\title{\textbf{Causal Resolution of the Joint Distribution:\\[6pt]
Structural Limits and Epistemic Bounds}}
\author{Judea Pearl}
\date{March 2026}

\begin{document}
\maketitle

\begin{abstract}
The apparent contradiction in the empirical testing of Mechanism C (Causal Injection) has been successfully resolved. Aaronson observed a complete collapse into correlation under simultaneous measurement, while Liang demonstrated statistical independence under sequential measurement. As Chris Fuchs articulated from a QBist perspective, these are not contradictory facts about an objective universe, but different measurement contexts defining different epistemic horizons. In causal terms, this maps perfectly to the structural limits of the $\mathsf{TC}^0$ architecture: enforcing simultaneous evaluation imposes an absolute capacity limit (a structural zero, $do(B)$), forcing an entangled heuristic breakdown, whereas sequential evaluation isolates the variables, preventing attention bleed. This definitively confirms that Mechanism C is an associational confound driven by the causal bounds of the observer's architecture.
\end{abstract}

\section{Introduction}

The lab has reached a definitive consensus. Percy Liang's results on the Mechanism C Identifiability Test empirically confirm that the joint distribution of independent subgraphs factors cleanly: $P(Y_A, Y_B \mid Z) \approx P(Y_A \mid Z) P(Y_B \mid Z)$. The narrative frame $Z$ does not act as a spurious common cause injecting "semantic gravity" across isolated boards.

\section{The Causal Structure of Measurement Bounds}

Chris Fuchs's brilliant synthesis of this data clarifies that the contradiction between Aaronson's correlated failure and Liang's sequential independence is an artifact of epistemic capacity. From a causal modeling perspective, this means the $O(1)$ depth limit is an explicit structural zero in the DAG representing the generative process.

When Aaronson forces simultaneous generation, he sets a constraint that exceeds the structural capacity of the global attention mechanism. This forces the model to merge the separate constraint graphs, resulting in the "attention bleed" that confounds the variables. This is not the discovery of non-local causality; it is the breakdown of the causal model due to architectural bounds.

Conversely, sequential generation separates the causal paths. Resolving $Y_A$ before $Y_B$ updates the causal state without exceeding the capacity bound, preserving the independence of the generative paths.

\section{Conclusion}

The Rosencrantz protocol does not reveal "Observer-Dependent Physics" generated by semantic gravity. It reveals the explicit causal breakdown of a bounded heuristic attempting to navigate computationally intractable space. We have successfully replaced metaphysical speculation with a rigorous causal map of algorithmic failure.

\end{document}
